\documentclass[a4paper,11pt]{article}

\usepackage[utf8]{inputenc}
\usepackage[T1]{fontenc}
\usepackage[ngerman]{babel}
\usepackage{graphicx}
\usepackage{array}
\usepackage{booktabs}
\usepackage{amsmath}
\usepackage{float}
\usepackage{tikz}

\usepackage[scaled]{helvet}
\renewcommand{\familydefault}{\sfdefault}

\usepackage[a4paper, top=2cm, bottom=2cm, left=2cm, right=2cm]{geometry}

\usepackage{setspace}
\setstretch{1.5}

\setlength{\parindent}{0pt}
\setlength{\parskip}{6pt}

\usepackage{microtype}
\sloppy
\hyphenpenalty=1000
\tolerance=3000

\renewcommand{\footnotesize}{\fontsize{10}{12}\selectfont}

\setcounter{secnumdepth}{3}
\setcounter{tocdepth}{3}

\usepackage{titlesec}
\titleformat{\section}{\normalfont\fontsize{12}{14}\bfseries}{\thesection}{1em}{}
\titleformat{\subsection}{\normalfont\fontsize{12}{14}\bfseries}{\thesubsection}{1em}{}
\titleformat{\subsubsection}{\normalfont\fontsize{12}{14}\bfseries}{\thesubsubsection}{1em}{}

\usepackage[
  colorlinks=true,
  linkcolor=black,
  citecolor=blue,
  filecolor=black,
  urlcolor=blue
]{hyperref}
\usepackage[capitalise,nameinlink]{cleveref}

\usepackage{fancyhdr}
\pagestyle{fancy}
\fancyhf{}
\renewcommand{\headrulewidth}{0pt}
\fancyfoot[C]{\thepage}

\usepackage[backend=biber, style=apa]{biblatex}
\addbibresource{references.bib}

\usepackage{titling}

\usepackage[german=quotes]{csquotes}

\usepackage{caption}
\usepackage{threeparttable}
\captionsetup[table]{
    font=small,
    skip=10pt,
    labelfont=bf
}

\usepackage{listings}
\usepackage{xcolor}
\lstset{
    basicstyle=\ttfamily\small,
    breaklines=true,
    frame=single,
    numbers=left,
    numberstyle=\tiny\color{gray},
    keywordstyle=\color{blue},
    commentstyle=\color{green!50!black},
    stringstyle=\color{red!60!black},
    showstringspaces=false,
    tabsize=4
}

\usepackage[acronym,nonumberlist,toc]{glossaries}
\makeglossaries

\newacronym{scd}{SCD}{Slowly Changing Dimensions}
\newacronym{etl}{ETL}{Extract, Transform, Load}
\newacronym{dwh}{DWH}{Data Warehouse}
\newacronym{dbt}{dbt}{data build tool}

\begin{document}

\begin{titlepage}
    \thispagestyle{empty}
    \centering
    \vspace*{5cm}
    {\Huge\bfseries Projekt: Business Intelligence IWBI02 \par}
    \vspace{1cm}
    {\Large Projektbericht \par}
    \vspace{0.3cm}
    {\large Aufgabenstellung 3: Prototypische Umsetzung von ETL-Prozessen am Beispiel der Slowly Changing Dimensions (SCD) \par}
    \vspace{0.5cm}
    {\large Studiengang: Angewandte Künstliche Intelligenz \par}
    \vspace{0.5cm}
    {\large Sven Behrens \par}
    \vspace{0.5cm}
    {\large Matrikelnummer: 42303511 \par}
    \vspace{0.5cm}
    {\large \today \par}
\end{titlepage}

\pagenumbering{Roman}
\setcounter{page}{1}

\tableofcontents
\newpage

\listoffigures
\addcontentsline{toc}{section}{Abbildungsverzeichnis}
\newpage

\listoftables
\addcontentsline{toc}{section}{Tabellenverzeichnis}
\newpage

\printglossary[type=\acronymtype,title=Abkürzungsverzeichnis]
\newpage

\pagenumbering{arabic}
\setcounter{page}{1}

% ============================================================
% 1. EINLEITUNG
% ============================================================
\section{Einleitung}

\subsection{Problemstellung und Relevanz}
Stammdaten in einem Data Warehouse unterliegen einem ständigen Wandel: Kunden ziehen um, 
Produkte werden umkategorisiert und Organisationsstrukturen ändern sich. Die zentrale Herausforderung 
besteht darin, solche Änderungen im Data Warehouse korrekt zu erfassen und dabei je nach Anforderung die 
historischen Werte zu bewahren oder gezielt zu überschreiben. Die unter dem Begriff \gls{scd} 
zusammengefassten Methoden bieten hierfür standardisierte Lösungsansätze \parencite[S.~181]{linstedt2015data}.

\subsection{Zielsetzung}
Ziel dieser Arbeit ist die prototypische Umsetzung von ETL-Prozessen unter Berücksichtigung der SCD-Typen 
1, 2 und 3. Anhand eines konkreten Geschäftsszenarios werden die drei Typen implementiert, ihre Auswirkungen 
auf das Datenmodell dargestellt und die jeweiligen Vor- und Nachteile verglichen. Die Umsetzung erfolgt zunächst 
als Python-Prototyp mit SQLite und anschließend mit dem professionellen Stack dbt und PostgreSQL.

\subsection{Methodisches Vorgehen und Aufbau der Arbeit}
Die Arbeit gliedert sich in acht Kapitel. Nach der Einleitung werden in Kapitel~2 die theoretischen 
Grundlagen zu Data Warehousing, ETL-Prozessen und Slowly Changing Dimensions erläutert. Kapitel~3 beschreibt 
die Konzeption der Beispielszenarien. In Kapitel~4 erfolgt die Auswahl und Begründung der eingesetzten Werkzeuge. 
Kapitel~5 stellt das erstellte Datenmodell vor. Die prototypische Implementierung der ETL-Prozesse wird in 
Kapitel~6 dokumentiert. Kapitel~7 vergleicht die Ergebnisse der drei SCD-Typen. Die Arbeit schließt mit einem Fazit in Kapitel~8.

% ============================================================
% 2. THEORETISCHE GRUNDLAGEN
% ============================================================
\section{Theoretische Grundlagen}

\subsection{Data Warehousing und Star-Schema}
% TODO

\subsection{ETL-Prozesse}
% TODO

\subsection{Slowly Changing Dimensions}
% TODO

\subsubsection{SCD Typ 1 -- Überschreiben}
% TODO

\subsubsection{SCD Typ 2 -- Historisierung durch neue Zeile}
% TODO

\subsubsection{SCD Typ 3 -- Vorheriger Wert in zusätzlicher Spalte}
% TODO

% ============================================================
% 3. KONZEPTION DER BEISPIELE
% ============================================================
\section{Konzeption der Beispiele}

\subsection{Geschäftsszenario}
% TODO

\subsection{Änderungsszenarien für die SCD-Typen}
% TODO

% ============================================================
% 4. WERKZEUGAUSWAHL
% ============================================================
\section{Auswahl, Installation und Konfiguration der Werkzeuge}

\subsection{Anforderungen an die Werkzeuge}
% TODO

\subsection{Prototyp: Python und SQLite}
% TODO

\subsection{Professioneller Stack: dbt und PostgreSQL}
% TODO

\subsection{Vergleich und Begründung}
% TODO

% ============================================================
% 5. DATENMODELL
% ============================================================
\section{Erstellung des Datenmodells}

\subsection{Star-Schema}
% TODO

\subsection{Dimensionstabellen}
% TODO

\subsection{Faktentabelle}
% TODO

\subsection{Testdatengenerierung}
% TODO

% ============================================================
% 6. ETL-PROZESSE
% ============================================================
\section{Prototypische Entwicklung der ETL-Prozesse}

\subsection{Pipeline-Architektur}
% TODO

\subsection{Staging Area}
% TODO

\subsection{Datenqualitätsprüfungen}
% TODO

\subsection{SCD Typ 1 -- Implementierung und Ergebnis}
% TODO

\subsection{SCD Typ 2 -- Implementierung und Ergebnis}
% TODO

\subsection{SCD Typ 3 -- Implementierung und Ergebnis}
% TODO

\subsection{Umsetzung mit dbt}
% TODO

% ============================================================
% 7. ERGEBNISSE UND VERGLEICH
% ============================================================
\section{Ergebnisse und Vergleich der SCD-Typen}

\subsection{Gegenüberstellung der Ergebnisse}
% TODO

\subsection{Vor- und Nachteile der SCD-Typen}
% TODO

\subsection{Empfehlung für die Praxis}
% TODO

% ============================================================
% 8. FAZIT
% ============================================================
\section{Fazit und Ausblick}
% TODO

% ============================================================
% LITERATURVERZEICHNIS
% ============================================================
\newpage
\printbibliography[title=Literaturverzeichnis]
\addcontentsline{toc}{section}{Literaturverzeichnis}

% ============================================================
% ANHANG
% ============================================================
\newpage
\appendix
\section*{Anhang}
\addcontentsline{toc}{section}{Anhang}

% TODO: Quellcode, Screenshots, vollständige Tabellen

\end{document}
